\chapter{Tropical sprue}
\index{Tropical sprue}

\section*{Definition and Overview}
Tropical sprue is a chronic, acquired, progressive diarrhoeal disease characterized by malabsorption of essential nutrients from the small intestine, most notably folic acid (folate) and vitamin B12 (cobalamin). It is endemic to specific tropical and subtropical regions of the world, particularly parts of the Caribbean, the Indian subcontinent, Southeast Asia, and Central and South America. The name itself reflects its geographic distribution and its primary clinical presentation: a malabsorptive state (``sprue,'' derived from the Dutch word \textit{spruw}, which historically described aphthous disease or thrush of the mouth) occurring in tropical climates.

Historically, the disease was a major cause of morbidity and mortality among military personnel, colonial administrators, and long-term travelers visiting endemic areas. While it has become less common in some regions due to improvements in sanitation, access to clean drinking water, and public health infrastructure, it remains a significant health concern for local populations and a vital diagnostic consideration for returning travelers presenting with chronic diarrhoea, nutritional deficiencies, and progressive weight loss.

The condition is histologically characterized by progressive structural damage to the mucosal lining of the small intestine. This damage typically begins in the distal duodenum and jejunum, later progressing to involve the ileum. The hallmark of the disease is a chronic inflammatory response that leads to blunting or flattening of the intestinal villi (villous atrophy), elongation of the intestinal crypts of Lieberkuhn (crypt hyperplasia), and infiltration of the lamina propria with chronic inflammatory cells, such as lymphocytes, plasma cells, and eosinophils. This structural damage severely impairs the mucosal surface area available for nutrient absorption, leading to progressive nutritional deficiencies, significant weight loss, wasting, and systemic complications such as megaloblastic anaemia.

Unlike celiac disease (gluten-sensitive enteropathy), which is a lifelong autoimmune reaction to dietary gluten, tropical sprue is an acquired condition that does not respond to gluten restriction but is highly responsive to long-term broad-spectrum antibiotic therapy and nutritional supplementation. Its exact primary trigger remains a subject of ongoing investigation, though a persistent bacterial infection or post-infectious dysbiosis is widely believed to be the root cause. Understanding tropical sprue is essential for healthcare providers, as early recognition and treatment can completely reverse the intestinal damage and prevent long-term systemic complications.

\section*{Epidemiology}
Tropical sprue is strictly localized to certain geographical regions within the tropics and subtropics, generally between 30 degrees north and 30 degrees south of the equator. The disease is endemic in the Indian subcontinent (particularly southern and northern rural areas of India, Pakistan, and Bangladesh), parts of Southeast Asia (such as Burma, Vietnam, Thailand, and Indonesia), the Caribbean islands (including Puerto Rico, Cuba, Haiti, and the Dominican Republic), and portions of Central and South America (such as Colombia and Venezuela). Curiously, despite sharing similar tropical climates, agricultural practices, and levels of sanitation, tropical sprue is rare or virtually non-existent in sub-Saharan Africa. This geographic disparity remains an unsolved mystery in medical literature.

The disease affects both local residents in endemic regions and visitors who travel to or reside in these areas for extended periods. Historically, epidemics of tropical sprue were documented among British troops stationed in India and American military personnel in East Asia and the Caribbean during the 20th century. Today, cases in non-endemic countries are almost exclusively diagnosed in immigrants, long-term expatriates, missionaries, and military personnel returning from endemic zones. Brief visits of less than a few weeks rarely result in tropical sprue; most affected travelers have resided in an endemic area for months or years before symptoms manifest.

In terms of demographics, the disease does not show a significant predilection for any specific age group, though it is more commonly diagnosed in adults between the ages of 20 and 50. There is no clear gender bias, with males and females affected in equal proportions. The incidence of tropical sprue has significantly declined over the past few decades, a trend largely attributed to widespread improvements in sanitation, the availability of clean drinking water, and the early use of broad-spectrum antibiotics for acute gastrointestinal infections. Despite this decline, it remains a common cause of chronic diarrhoea and malabsorption in low-resource rural settings within endemic areas, where it contributes to chronic malnutrition and public health burdens.

\section*{Aetiology and Pathophysiology}
The precise cause of tropical sprue remains partially elusive, but the consensus view is that it represents a chronic, post-infectious malabsorptive state caused by persistent bacterial overgrowth or colonization of the small intestine by toxigenic bacteria. The pathogenetic process involves a complex interaction between environmental pathogens, host mucosal immunity, and nutritional depletion.

Key factors in the aetiology and pathophysiology include:
\begin{itemize}
    \item \textbf{Persistent Bacterial Colonization:} Unlike the healthy small intestine, which contains relatively few bacteria, the small intestine of individuals with tropical sprue is heavily colonized by coliform bacteria. The most commonly isolated organisms include toxigenic strains of \textit{Klebsiella pneumoniae}, \textit{Escherichia coli}, and \textit{Enterobacter cloacae}. These bacteria produce enterotoxins and ethanol, which injure the enterocytes (intestinal epithelial cells) and impair mucosal enzymatic activity.
    \item \textbf{Small Intestinal Bacterial Overgrowth (SIBO):} Persistent colonization leads to a state of chronic bacterial overgrowth. These bacteria deconjugate bile acids, which impairs micelles formation and leads to the malabsorption of dietary lipids. The bacteria also compete with the host for essential nutrients, particularly vitamin B12 (cobalamin), contributing to systemic deficiencies.
    \item \textbf{Intestinal Mucosal Injury:} The chronic presence of pathogenic bacteria and their metabolic byproducts triggers a sustained inflammatory response. Histologically, this leads to mucosal damage characterized by varying degrees of villous atrophy (ranging from mild blunting to total flattening of the villi), crypt hyperplasia, and a marked increase in intraepithelial lymphocytes and chronic inflammatory cells (plasma cells and lymphocytes) in the lamina propria. The loss of functional villous surface area severely diminishes the intestine's capacity to absorb water, electrolytes, and macronutrients.
    \item \textbf{Nutritional Deficiencies and Mucosal Regeneration Impairment:} The malabsorption of folate and vitamin B12 creates a destructive feedback loop. Folate is crucial for rapid cell division, including the constant renewal of the intestinal epithelium. As folate stores are depleted (exacerbated by malabsorption), mucosal regeneration is impaired, leading to further mucosal atrophy and worsening malabsorption.
    \item \textbf{Gastrointestinal Motility Alterations:} The structural damage and chronic inflammation disrupt normal small bowel motility. The migrating motor complex (MMC), which acts as a physiological ``sweeper'' to clear bacteria from the small intestine, is often impaired. This hypomotility further perpetuates bacterial overgrowth, locking the patient into a chronic cycle of infection, inflammation, and malabsorption.
\end{itemize}

\section*{Clinical Features}
The clinical presentation of tropical sprue typically evolves through distinct phases, beginning with acute gastrointestinal symptoms and progressing to a chronic state characterized by severe nutritional depletion and systemic manifestations. The onset can be insidious or may follow an initial episode of acute infective diarrhoea (such as traveler's diarrhoea or food poisoning) that fails to resolve.

The signs and symptoms of tropical sprue include:
\begin{itemize}
    \item \textbf{Chronic Diarrhoea:} The primary and most common symptom is persistent, watery, or loose stools. As malabsorption worsens, the diarrhoea often takes on the characteristics of steatorrhea, where the stools become voluminous, pale, greasy, foul-smelling, and difficult to flush due to high fat content.
    \item \textbf{Abdominal Symptoms:} Patients frequently complain of abdominal distension, bloating, flatulence, and a persistent sensation of fullness. Cramping or mild abdominal pain is common and is often exacerbated by eating.
    \item \textbf{Significant Weight Loss and Wasting:} Due to the malabsorption of fats, proteins, and carbohydrates, patients experience progressive weight loss. In chronic cases, this can lead to muscle wasting (cachexia), severe lethargy, and profound physical weakness.
    \item \textbf{Megaloblastic Anaemia Symptoms:} A hallmark feature resulting from the combined malabsorption of folate and vitamin B12 is megaloblastic anaemia. This manifests as severe fatigue, generalized weakness, exertional dyspnoea (shortness of breath), palpitations, and a pale or slightly icteric (jaundiced) complexion.
    \item \textbf{Mucosal and Epithelial Changes:} Nutritional deficiencies lead to visible epithelial changes. Glossitis (a painful, smooth, beefy-red tongue), angular cheilitis (painful cracks at the corners of the mouth), and recurrent aphthous ulcers are common. Hyperpigmentation of the skin, dry skin, and brittle nails may also occur.
    \item \textbf{Neurological Manifestations:} Prolonged vitamin B12 deficiency can lead to neurological complications. Patients may experience peripheral neuropathy, characterized by paresthesia (tingling or numbness) in the hands and feet, loss of vibratory and position sense, ataxia (difficulty walking), and cognitive changes such as irritability or mild memory impairment.
    \item \textbf{Systemic Nutritional Deficiencies:} Malabsorption of vitamin D and calcium can lead to bone pain, muscle cramps, and tetany. Vitamin K deficiency can cause easy bruising or bleeding tendencies (petechiae and ecchymoses). Edema of the lower extremities may occur due to hypoalbuminaemia resulting from protein-losing enteropathy and poor protein absorption.
\end{itemize}

\section*{Differential Diagnosis}
Given that chronic diarrhoea, malabsorption, and weight loss are features of many gastrointestinal conditions, the differential diagnosis for tropical sprue is broad. It is crucial to systematically rule out other causes, particularly those requiring different therapeutic interventions.

Key conditions to consider in the differential diagnosis include:
\begin{itemize}
    \item \textbf{Celiac Disease:} This is an autoimmune enteropathy triggered by the ingestion of gluten in genetically susceptible individuals. While it presents with identical histological features (villous atrophy, crypt hyperplasia) and clinical symptoms, celiac disease is lifelong, does not respond to antibiotics, and is managed with a strict gluten-free diet. It is associated with specific serological markers, such as anti-tissue transglutaminase (tTG) antibodies, which are absent in tropical sprue.
    \item \textbf{Crohn's Disease:} A chronic inflammatory bowel disease that can affect any part of the gastrointestinal tract, including the small intestine. Unlike tropical sprue, Crohn's disease is characterized by transmural inflammation, skip lesions, granulomas on biopsy, and potential complications like fistulae and strictures.
    \item \textbf{Chronic Giardiasis:} Infection with the protozoan parasite \textit{Giardia duodenalis} (also known as \textit{Giardia lamblia}) can cause chronic diarrhoea, bloating, steatorrhea, and weight loss, mimicking tropical sprue. It is diagnosed by identifying cysts or trophozoites in stool samples or via duodenal fluid aspiration.
    \item \textbf{Small Intestinal Bacterial Overgrowth (SIBO):} While SIBO is a component of tropical sprue's pathophysiology, SIBO can occur independently due to anatomical abnormalities (e.g., blind loops, diverticula), motility disorders (e.g., scleroderma, diabetic enteropathy), or chronic proton pump inhibitor use.
    \item \textbf{Whipple's Disease:} A rare systemic infectious disease caused by the bacterium \textit{Tropheryma whipplei}. It presents with malabsorption, weight loss, joint pain, and neurological signs. Small bowel biopsy reveals periodic acid-Schiff (PAS)-positive macrophages containing the bacilli, distinguishing it from tropical sprue.
    \item \textbf{HIV-Associated Enteropathy:} Chronic diarrhoea and malabsorption are common in individuals with advanced HIV/AIDS, often due to opportunistic infections (e.g., \textit{Cryptosporidium}, \textit{Microsporidia}, \textit{Mycobacterium avium} complex) or the direct cytopathic effect of the virus on the enterocytes.
    \item \textbf{Intestinal Tuberculosis:} A common diagnostic consideration in developing countries where tuberculosis is endemic. It can present with chronic diarrhoea, malabsorption, and weight loss, but typically features caseating granulomas on biopsy and lung involvement on chest imaging.
    \item \textbf{Pancreatic Insufficiency:} Chronic pancreatitis or cystic fibrosis can cause malabsorption and steatorrhea due to a lack of digestive enzymes. However, the small bowel mucosa remains normal, and fecal elastase levels are typically low.
\end{itemize}

\section*{When to Seek Medical Care}
Persistent gastrointestinal symptoms should never be ignored, particularly after returning from travel to tropical or subtropical regions. While many acute diarrhoeal illnesses resolve spontaneously, chronic symptoms require a thorough medical evaluation to prevent severe nutritional depletion and long-term complications.

A medical consultation should be scheduled if any of the following symptoms occur:
\begin{itemize}
    \item Diarrhoea that persists for more than two weeks, especially if accompanied by foul-smelling, greasy, or unusually pale stools.
    \item Unexplained and progressive weight loss despite a normal or increased appetite.
    \item Chronic abdominal bloating, distension, or persistent cramping pain after meals.
    \item Extreme fatigue, generalized weakness, or a pale appearance that does not improve with rest.
    \item Soreness or redness of the tongue, cracks at the corners of the mouth, or recurrent mouth ulcers.
    \item Numbness, tingling, or a ``pins and needles'' sensation in the hands or feet, or difficulty maintaining balance when walking.
\end{itemize}

Immediate emergency medical care must be sought if any of the following life-threatening symptoms are present. In many countries, call local emergency services; in the United States or Canada, call 911; in the United Kingdom, call 999; or contact the local emergency services immediately:
\begin{itemize}
    \item Signs of severe, life-threatening dehydration, such as extreme thirst, dry mouth, little to no urination, dark urine, confusion, dizziness, or lightheadedness upon standing.
    \item High, persistent fever (above 38.5 degrees Celsius or 101.3 degrees Fahrenheit) accompanied by severe abdominal pain.
    \item Inability to keep fluids down due to persistent vomiting.
    \item Passage of bloody stools or black, tarry stools (melaena).
    \item Shortness of breath, chest pain, or rapid heart rate, which may indicate severe anaemia or cardiovascular compromise.
\end{itemize}

\section*{Diagnosis and Investigations}
The diagnosis of tropical sprue is established through a combination of clinical history (confirming exposure to an endemic region), laboratory evidence of malabsorption and nutritional deficiencies, characteristic histological findings on small bowel biopsy, and the exclusion of other diseases. Importantly, a positive response to treatment with antibiotics and folic acid confirms the diagnosis.

Diagnostic investigations typically include:
\begin{itemize}
    \item \textbf{Blood Tests:} A complete blood count (CBC) often reveals megaloblastic anaemia, characterized by a high mean corpuscular volume (MCV), low haemoglobin, and hypersegmented neutrophils. Serum levels of folate and vitamin B12 are characteristically low. Comprehensive metabolic panels may show hypoalbuminaemia (due to malabsorption and protein-losing enteropathy) and electrolyte imbalances (low potassium, magnesium, and calcium).
    \item \textbf{Stool Studies:} Stool examination is critical to rule out parasitic, bacterial, or viral pathogens. Stool cultures, microscopy for ova and parasites, and antigen testing (e.g., for \textit{Giardia}) are performed. Stool fat analysis (such as a 72-hour fecal fat collection) can quantify the degree of steatorrhea, and fecal elastase is checked to rule out exocrine pancreatic insufficiency.
    \item \textbf{Malabsorption Tests:} The d-xylose absorption test is historically used to assess carbohydrate absorption. Since d-xylose is absorbed in the proximal small intestine, low urinary excretion of d-xylose suggests mucosal damage in the duodenum or jejunum, supporting a diagnosis of tropical sprue.
    \item \textbf{Upper Gastrointestinal Endoscopy and Biopsy:} This is the gold standard diagnostic tool. Endoscopy allows direct visualization of the duodenum, which may show loss of mucosal folds, scalloping, or a mosaic pattern. Multiple biopsy samples are obtained from the second and third parts of the duodenum.
    \item \textbf{Histopathology:} Microscopic examination of the tissue biopsy samples reveals characteristic, though non-specific, abnormalities:
        \begin{itemize}
            \item Blunting or complete atrophy of the intestinal villi.
            \item Elongation of the crypts of Lieberkuhn (crypt hyperplasia).
            \item An increased number of intraepithelial lymphocytes (IELs).
            \item Chronic inflammatory cell infiltration (lymphocytes, plasma cells, and eosinophils) within the lamina propria.
        \end{itemize}
        The severity of these changes typically correlates with the duration and clinical severity of the disease.
\end{itemize}

\section*{Management}
The management of tropical sprue is highly effective and focuses on two primary goals: eradicating the underlying bacterial overgrowth and correcting the secondary nutritional deficiencies. With adherence to therapy, the vast majority of patients achieve complete clinical and histological recovery.

Treatment strategies include:
\begin{itemize}
    \item \textbf{Antibiotic Therapy:} A prolonged course of antibiotics is the cornerstone of treatment to eliminate toxigenic bacterial colonization.
        \begin{itemize}
            \item \textbf{Tetracycline:} The traditional first-line antibiotic is tetracycline, administered at a dose of 250 mg four times daily (or 500 mg twice daily) for a period of 3 to 6 months. For patients with long-standing disease or severe mucosal damage, treatment may be extended to 12 months.
            \item \textbf{Alternative Antibiotics:} In patients for whom tetracycline is contraindicated (such as pregnant women and children under the age of 8, due to the risk of permanent tooth discoloration), alternative options include doxycycline (100 mg twice daily), sulfamethoxazole-trimethoprim (co-trimoxazole), or broad-spectrum agents like rifaximin, which remains largely confined to the gastrointestinal tract.
        \end{itemize}
    \item \textbf{Folic Acid Supplementation:} Folic acid is administered orally at a dose of 5 mg daily. Folic acid not only corrects systemic folate deficiency but also stimulates the rapid regeneration of the small intestinal mucosa, leading to a prompt improvement in symptoms. In mild or early cases, folic acid supplementation alone can lead to clinical improvement, though it should be combined with antibiotics to ensure complete eradication of the disease.
    \item \textbf{Vitamin B12 Supplementation:} For patients with documented vitamin B12 deficiency or severe megaloblastic anaemia, intramuscular injections of cobalamin (typically 1000 micrograms weekly for several weeks, followed by monthly injections) are administered. Oral supplementation may be insufficient initially due to the severe mucosal damage in the ileum, where vitamin B12 is absorbed.
    \item \textbf{Nutritional Support and Rehydration:} Oral or intravenous rehydration is initiated in patients presenting with dehydration or electrolyte imbalances. A high-protein, low-fat diet may be recommended initially to reduce steatorrhea and support tissue repair. Supplemental fat-soluble vitamins (A, D, E, and K), calcium, and iron are prescribed as indicated by laboratory testing.
    \item \textbf{Monitoring and Follow-up:} Patients require regular clinical monitoring to assess response to therapy. Resolution of diarrhoea, weight gain, and normalization of haematological parameters are expected within weeks of starting treatment. A repeat small bowel biopsy is occasionally performed in refractory cases to confirm mucosal healing.
\end{itemize}

\section*{Complications}
When diagnosed early and treated appropriately, tropical sprue rarely leads to long-term complications. However, if the condition remains untreated or if therapy is delayed, the chronic state of malabsorption and inflammation can result in significant morbidity and life-threatening complications.

Potential complications of untreated or severe tropical sprue include:
\begin{itemize}
    \item \textbf{Severe Megaloblastic Anaemia:} Prolonged deficiency of folate and vitamin B12 impairs red blood cell production, leading to severe anaemia. This can cause high-output heart failure, particularly in elderly patients or those with pre-existing cardiovascular disease.
    \item \textbf{Irreversible Neurological Damage:} Long-standing vitamin B12 deficiency leads to subacute combined degeneration of the spinal cord (SCD). This condition causes progressive demyelination of the posterior and lateral columns of the spinal cord, resulting in permanent sensory loss, spastic paraparesis, ataxia, and cognitive decline if not treated promptly.
    \item \textbf{Severe Malnutrition and Cachexia:} The inability to absorb macronutrients leads to profound weight loss, muscle wasting, and loss of subcutaneous fat. This state of severe malnutrition impairs immune function, rendering the patient highly susceptible to secondary opportunistic infections.
    \item \textbf{Electrolyte and Mineral Imbalances:} Chronic diarrhoea and malabsorption lead to depletion of potassium (hypokalaemia), magnesium (hypomagnesaemia), and calcium (hypocalcaemia). These imbalances can cause life-threatening cardiac arrhythmias, muscle tetany, seizures, and metabolic bone disease (osteomalacia or osteoporosis).
    \item \textbf{Bleeding Diathesis:} Malabsorption of the fat-soluble vitamin K impairs the synthesis of clotting factors in the liver. This can lead to easy bruising, mucosal bleeding, hematuria, and in severe cases, life-threatening internal haemorrhage.
    \item \textbf{Growth Retardation in Children:} In the rare cases where tropical sprue affects children, chronic malabsorption can lead to severe growth failure, delayed puberty, and developmental delays.
\end{itemize}

\section*{Prognosis and Outcomes}
The prognosis for individuals with tropical sprue is excellent, provided they receive early diagnosis and appropriate, complete treatment. The vast majority of patients experience a rapid and dramatic response to therapy. Diarrhoea and abdominal symptoms typically begin to improve within the first week of starting folic acid and antibiotic therapy, and nutritional parameters, including weight gain and haematological values, return to normal within a few months.

Histological recovery of the small bowel mucosa is slower but generally complete. In most cases, repeat biopsies demonstrate normal mucosal architecture within 3 to 6 months of completing therapy.

Factors that influence the prognosis and long-term outcomes include:
\begin{itemize}
    \item \textbf{Compliance with Therapy:} Adherence to the full course of antibiotics and nutritional supplements is the single most critical factor in achieving a cure. Premature discontinuation of antibiotics can lead to recurrence of bacterial overgrowth and relapse of symptoms.
    \item \textbf{Duration of Symptoms Before Treatment:} Patients who have suffered from untreated tropical sprue for many years may take longer to recover, and those who have developed neurological complications from vitamin B12 deficiency may be left with residual, irreversible neurological deficits.
    \item \textbf{Return to Endemic Areas:} For travelers who return to non-endemic countries, recurrence is virtually unknown once the disease is cured. However, for local residents in endemic regions, there is a risk of reinfection and relapse, particularly if they continue to live in environments with poor sanitation and contaminated water supplies. In these populations, ongoing public health interventions are necessary to prevent recurrence.
\end{itemize}

\section*{Prevention and Self-Care}
Because the exact primary trigger of tropical sprue is linked to contamination and bacterial exposure in endemic regions, prevention strategies are closely aligned with those used to prevent other waterborne and foodborne diarrhoeal diseases. For long-term residents and travelers to endemic areas, practicing meticulous food and water hygiene is the most effective way to reduce risk.

Recommended prevention and self-care measures include:
\begin{itemize}
    \item \textbf{Safe Water Practices:} Avoid drinking untreated tap water in endemic areas. Consume only bottled water from sealed, reputable brands, or water that has been boiled for at least one minute (or three minutes at higher altitudes). Avoid ice cubes in beverages, as they are often made from local tap water.
    \item \textbf{Food Safety Hygiene:} Practice the classic travel rule: ``Cook it, peel it, forget it.'' Eat food that is freshly cooked and served hot. Avoid raw fruits and vegetables unless they can be peeled by hand after thorough washing with safe water. Avoid raw or undercooked meats, seafood, and unpasteurized dairy products. Avoid food from street vendors where hygiene standards may be low.
    \item \textbf{Personal Hand Hygiene:} Wash hands frequently with soap and clean running water, especially before preparing food, eating, and after using the bathroom. If clean water and soap are unavailable, use an alcohol-based hand sanitizer containing at least 60\% alcohol.
    \item \textbf{Early Treatment of Acute Diarrhoea:} Promptly treat episodes of acute travelers' diarrhoea or gastroenteritis with hydration and, if clinically indicated, short courses of antibiotics under medical supervision. This may prevent the establishment of chronic bacterial colonization that can lead to tropical sprue.
    \item \textbf{Self-Monitoring and Nutritional Support:} During recovery, adhere strictly to the prescribed antibiotic and supplementation regimen. Monitor body weight and symptoms, and report any signs of recurrence or worsening to a doctor. Consume a balanced, nutrient-dense diet rich in proteins and vitamins to replenish depleted body stores.
\end{itemize}

\section*{Sources and Further Reading}
For reliable, evidence-based information and further reading on tropical sprue, please refer to the following clinical and educational resources:
\begin{itemize}
    \item \textbf{World Health Organization (WHO):} International travel and health guidelines, providing updates on sanitation, infectious diseases, and geographic health risks. Website: \url{https://www.who.int}
    \item \textbf{Centers for Disease Control and Prevention (CDC):} The CDC Yellow Book: Health Information for International Travel, which offers comprehensive advice on food and water precautions and endemic gastrointestinal disorders. Website: \url{https://wwwnc.cdc.gov/travel}
    \item \textbf{National Institute of Diabetes and Digestive and Kidney Diseases (NIDDK):} Part of the U.S. National Institutes of Health, providing patient-friendly and professional resources on malabsorption syndromes. Website: \url{https://www.niddk.nih.gov}
    \item \textbf{British Society of Gastroenterology (BSG):} Clinical guidelines and educational resources on chronic diarrhoea and small bowel disorders. Website: \url{https://www.bsg.org.uk}
    \item \textbf{Gastroenterological Society of Australia (GESA):} Fact sheets and clinical practice resources on malabsorptive diseases and gut health. Website: \url{https://www.gesa.org.au}
    \item \textbf{The Lancet Infectious Diseases / Gastroenterology \& Hepatology:} Peer-reviewed medical journals featuring up-to-date research and review articles on tropical enteropathies and sprue. Website: \url{https://www.thelancet.com}
\end{itemize}